% ============================================================================
\chapter{Installer ton environnement de travail}
\label{ch:environnement}
% ============================================================================

\begin{objectif}
\begin{itemize}
  \item Disposer d'un poste de d\'eveloppement complet pour les trois langages, gratuit.
  \item Comprendre ce qu'est un runtime, un compilateur, un serveur web, un conteneur.
  \item Installer et v\'erifier chaque outil, avec les liens officiels de t\'el\'echargement.
  \item Savoir r\'eparer une installation cass\'ee (le cas le plus fr\'equent d'abandon).
\end{itemize}
\end{objectif}

\section{Le mat\'eriel n\'ecessaire}

Un ordinateur de 2015 ou plus r\'ecent suffit largement.

\begin{tabz}{lX}
  \headrow \thd{Composant} \& \thd{Minimum vital} \\
  \midrule
  CPU \& 2 c\oe{}urs / 4 threads (i5 de 6\textsuperscript{e} g\'en\'eration, Ryzen~3, Apple~M1).\\
  RAM \& 8~Go pour suivre le livre. 16~Go si tu ouvres Docker + un IDE Java en m\^eme temps.\\
  Disque \& 40~Go libres (une distribution \LaTeX{} + JDK + SDK Android p\`esent lourd).\\
  \'Ecran \& Full~HD. Deux \'ecrans changent la vie : documentation d'un c\^ot\'e, \'editeur de l'autre.\\
  Internet \& Indispensable : paquets, d\'ependances, documentations.\\
\end{tabz}

Un ordinateur portable d'entr\'ee de gamme r\'ecent, Linux install\'e dessus, fait un poste de
d\'eveloppement excellent pour ces trois langages. Aucun besoin de machine de guerre.

\section{Vue d'ensemble : qu'allons-nous installer ?}

\begin{tabz}{llX}
  \headrow \thd{Outil} \& \thd{R\^ole} \& \thd{Officiel} \\
  \midrule
  \'Editeur / IDE \& L\`a o\`u tu \'ecris le code. \& \link{VS~Code}{https://code.visualstudio.com/} --- \link{VSCodium (sans telemetry)}{https://vscodium.com/} --- \link{Cursor}{https://www.cursor.com/} \\
  Node.js \& Runtime JavaScript hors navigateur + npm. \& \link{nodejs.org}{https://nodejs.org/fr/download} \\
  PHP \& Runtime + serveur int\'egr\'e + Composer. \& \link{php.net}{https://www.php.net/downloads.php} \\
  JDK \& Compile et ex\'ecute Java (JVM). \& \link{Adoptium Temurin}{https://adoptium.net/fr/} \\
  Git \& Historique de ton code. \& \link{git-scm.com}{https://git-scm.com/downloads} \\
  Docker \& Fait tourner MySQL, PostgreSQL, Redis sans les installer \og  pour de  \fg{}. \& \link{docker.com}{https://www.docker.com/products/docker-desktop/} \\
  MySQL / MariaDB \& Base de donn\'ees de tous les exemples PHP. \& \link{mysql.com}{https://dev.mysql.com/downloads/mysql/} \\
  PostgreSQL \& Base de donn\'ees des exemples Java. \& \link{postgresql.org}{https://www.postgresql.org/download/} \\
  Terminal \& Bash, Zsh, PowerShell, Windows~Terminal. \& \link{Windows~Terminal}{https://apps.microsoft.com/detail/9N0DX20HK701} \\
  Navigateur \& Chrome/Edge (outils de dev) + Firefox. \& \link{Chrome}{https://www.google.com/chrome/} \\
\end{tabz}

\subsection{Trois mots \`a comprendre tout de suite}

\begin{description}
  \item[Interpr\'et\'e vs compil\'e.] \langjs{} et \langphp{} sont \emph{ex\'ecut\'es} par un moteur qui lit
        le fichier source. \langjava{} est \emph{compil\'e} en bytecode (\fichier{.class}) ex\'ecut\'e
        par la JVM. Cons\'equence pratique : en JS/PHP tu modifies et tu rafra\^ichis ; en Java,
        tu recompiles (ton IDE ou Maven le fait pour toi).
  \item[Client vs serveur.] Le JavaScript de page s'ex\'ecute dans le navigateur de l'utilisateur
        (client). PHP, Node.js et Java s'ex\'ecutent sur une machine qui r\'epond \`a des requ\^etes
        (serveur). Un m\^eme langage peut faire les deux : c'est le cas de JavaScript.
  \item[Runtime vs SDK.] Le runtime sait \emph{ex\'ecuter}. Le SDK (JDK, PHP complet, Node) apporte
        en plus les outils pour \emph{d\'evelopper} : compilateur, gestionnaire de paquets,
        biblioth\`eque standard, utilitaires.
\end{description}

\section{Linux (Debian, Ubuntu, Fedora, Arch)}

C'est la plateforme la plus confortable pour ce livre : tout est \`a un \cmd{apt install}.

\etape{Node.js via nvm (la bonne m\'ethode)}
Ne jamais installer Node avec \cmd{apt} : la version du d\'ep\^ot est vieille de plusieurs ann\'ees.
On utilise \textbf{nvm} (Node Version Manager), qui permet d'installer, changer et r\'eparer
de versions en une commande.\src{https://github.com/nvm-sh/nvm\#installing-and-updating}

\begin{terminal}
# installer nvm (la commande officielle, verifie la version sur le depot)
curl -o- https://raw.githubusercontent.com/nvm-sh/nvm/v0.40.1/install.sh | bash
# relance le terminal, puis :
nvm install --lts          # installe la derniere version LTS
nvm use --lts
nvm alias default lts/*    # par defaut a chaque nouveau terminal
node -v && npm -v
\end{terminal}

Note : si tu g\`eres plusieurs langages (Node, PHP, Java, Python, Go), un gestionnaire de versions
unique est plus agr\'eable~: \textbf{mise} (\link{mise.jdx.dev}{https://mise.jdx.dev/}) ou
\textbf{asdf}. \link{mise}{https://github.com/jdx/mise} est \'ecrit en Rust, tr\`es rapide, compatible
avec les fichiers \fichier{.tool-versions} d'asdf.

\etape{PHP}
\begin{terminal}
# Debian/Ubuntu : le depot fournit un PHP recent + les extensions usuelles
sudo apt update
sudo apt install -y php-cli php-fpm php-mysql php-pgsql php-curl php-mbstring \
  php-xml php-zip php-gd php-intl php-bcmath php-sqlite3
php -v && php -m      # -m liste les extensions chargees

# Fedora
sudo dnf install php-cli php-fpm php-mysqlnd php-pgsql php-xml php-mbstring php-intl
# Arch
sudo pacman -S php php-pgsql php-sqlite
\end{terminal}

Sous Debian, la version du d\'ep\^ot est vieille. Pour un PHP r\'ecent~:
\link{deb.sury.org (paquets ONDrej sur mesure)}{https://deb.sury.org/}
\src{https://github.com/oerdnj/deb.sury.org}.

\etape{Java}
\begin{terminal}
# Temurin = OpenJDK certifie par l'Adoptium (Eclipse Foundation), gratuit et sans compte
sudo mkdir -m 0755 -p /etc/apt/keyrings
curl -fsSL https://packages.adoptium.net/artifactory/api/gpg/key/public \
  | sudo tee /etc/apt/keyrings/adoptium.asc >/dev/null
echo "deb [signed-by=/etc/apt/keyrings/adoptium.asc] https://packages.adoptium.net/artifactory/deb $(lsb_release -cs) main" \
  | sudo tee /etc/apt/sources.list.d/adoptium.list
sudo apt update && sudo apt install -y temurin-21-jdk
java -version && javac -version
echo 'export JAVA_HOME=/usr/lib/jvm/temurin-21-jdk-amd64' >> ~/.bashrc
\end{terminal}

Sous Arch~: \cmd{sudo pacman -S jdk21-openjdk}. Sous Fedora~:
\cmd{sudo dnf install java-21-temurin-devel || sudo dnf install java-21-openjdk-devel}.

\etape{Le reste}
\begin{terminal}
sudo apt install -y git curl wget unzip zip build-essential jq tree htop
# MySQL/MariaDB (ou laisse Docker le faire, section plus bas)
sudo apt install -y mariadb-server mariadb-client
# PostgreSQL
sudo apt install -y postgresql postgresql-contrib
# Docker (via le depot officiel, pas celui du distributeur)
curl -fsSL https://get.docker.com | sudo sh
sudo usermod -aG docker $USER   # deconnecte-toi puis reconnecte-toi
newgrp docker
docker run --rm hello-world
\end{terminal}

\section{macOS}

Installe \link{Homebrew}{https://brew.sh/fr/} : c'est le gestionnaire de paquets non officiel
mais universellement utilis\'e.

\begin{terminal}
/bin/bash -c "$(curl -fsSL https://raw.githubusercontent.com/Homebrew/install/HEAD/install.sh)"

# Node
brew install nvm && mkdir -p ~/.nvm
echo 'export NVM_DIR="$HOME/.nvm"' >> ~/.zprofile
echo '[ -s "$(brew --prefix nvm)/nvm.sh" ] && \. "$(brew --prefix nvm)/nvm.sh"' >> ~/.zprofile
source ~/.zprofile && nvm install --lts

# PHP : remplace le PHP systeme (vieux et bride) par une version brew complete
brew install php composer
brew link --overwrite php
echo 'export PATH="$(brew --prefix php)/bin:$PATH"' >> ~/.zprofile

# Java
brew install --cask temurin@21     # ou : brew install openjdk@21
/usr/libexec/java_home -V           # liste les JDK installes
echo 'export JAVA_HOME=$(/usr/libexec/java_home -v 21)' >> ~/.zprofile

# Les outils
brew install git tree jq wget
brew install --cask docker colortail iterm2
brew install --cask mariadb lnav || brew install mariadb
brew install --cask postgresql@16
\end{terminal}

\begin{note}
Sur Apple~Silicon, Homebrew installe dans \fichier{/opt/homebrew}. Si un outil est introuvable,
c'est presque toujours un probl\`eme de \verb|PATH|~: v\'erifie avec \cmd{echo $PATH} et
\cmd{which php}.
\end{note}

\section{Windows}

Deux strat\'egies, cumulables. Je recommande la~B pour le PHP/Java et la~A pour commencer
rapidement.

\etape{Strat\'egie A --- natif (le plus simple au d\'ebut)}
\begin{enumerate}
  \item \textbf{Node.js}~: installeur \link{nodejs.org/fr/download}{https://nodejs.org/fr/download}
        (LTS, .msi). coche \og ~Add to PATH~ \fg{}.
  \item \textbf{PHP}~: \link{windows.php.net/download}{https://windows.php.net/download/}
        --- ZIP \og ~VS16 x64 Thread Safe~ \fg{}, d\'ezippe dans \fichier{C:\php}, ajoute
        \fichier{C:\php} au \verb|PATH| (\emph{Param\`etres \$\rightarrow\$ Syst\`eme \$\rightarrow\$ \`A propos \$\rightarrow\$ Param\`etres
        syst\`eme avanc\'es \$\rightarrow\$ Variables d'environnement}), puis copie \fichier{php.ini-development}
        vers \fichier{php.ini} et d\'ecommentre les lignes
        \verb|extension=curl|, \verb|extension=mbstring|, \verb|extension=pdo\_mysql|,
        \verb|extension=mysqli|, \verb|extension=openssl|, \verb|extension=fileinfo|.
  \item \textbf{Java}~: \link{Adoptium Temurin 21 (Windows x64 MSI)}{https://adoptium.net/temurin/releases/?os=windows} ---
        coche \og ~Set JAVA\_HOME~ \fg{} et \og ~Add to PATH~ \fg{} pendant l'installation.
  \item \textbf{Git for Windows}~: \link{git-scm.com/download/win}{https://git-scm.com/download/win}
        (il fournit aussi Git~Bash, un vrai terminal Linux).
  \item \textbf{Windows~Terminal}~: \link{Microsoft~Store}{https://apps.microsoft.com/detail/9N0DX20HK701}.
  \item \textbf{Scoop} ou \textbf{winget} pour tout le reste~:
\end{enumerate}

\begin{terminal}
winget install --id Git.Git -e
winget install OpenJS.NodeJS.LTS
winget install PHP.PHP.8.3
winget install EclipseAdoptium.Temurin.21.JDK
winget install Docker.DockerDesktop
winget install Microsoft.VisualStudioCode
scoop bucket add extras; scoop install jq tree wget
\end{terminal}

\etape{Strat\'egie B --- WSL2 (un vrai Linux dans Windows, la voie royale)}
Tu gardes VS~Code sous Windows, mais tout le code et tous les outils vivent dans Ubuntu.
C'est l'environnement le plus proche de celui d'un serveur de production.

\begin{terminal}
# PowerShell en administrateur
wsl --install -d Ubuntu-24.04
wsl --update
wsl --set-default-version 2
# puis redemarre, ouvre "Ubuntu", cree ton utilisateur, et applique la section Linux ci-dessus
code .        # depuis le terminal WSL : ouvre VS Code sur le dossier distant
\end{terminal}

Installe l'extension \link{Remote - WSL}{https://marketplace.visualstudio.com/items?itemName=ms-vscode-remote.remote-wsl}
et \link{Dev~Containers}{https://marketplace.visualstudio.com/items?itemName=ms-vscode-remote.remote-containers}.
V\'erifie que ton projet est dans \verb|\\wsl\$\Ubuntu\home\toi\projets| et \emph{jamais} dans
\fichier{C:\textbackslash Users} : les performances d'I/O entre les deux mondes sont d\'esastreuses.

\begin{piege}
Sous Windows, les noms de fichiers sont insensibles \`a la casse, mais pas sur un serveur Linux.
Un \verb|require("./utils/Helper")| qui marche sur ton PC peut planter en production.
Convention~: noms de fichiers en \texttt{kebab-case} ou \texttt{snake\_case}, et imports
exactement identiques \`a la casse du fichier.
\end{piege}

\section{\'Editeur ou IDE : le choix}

\subsection{Option 1 --- un seul \'editeur pour les trois langages : VS~Code}

\begin{terminal}
# VS Code
https://code.visualstudio.com/                # Windows / macOS / Linux (.deb, .rpm, tar.gz)
# Variante sans telemetry : VSCodium
https://vscodium.com/
\end{terminal}

Extensions \`a installer (\keys{Ctrl+Shift+X}, puis cherche le nom~; le compte d'installations est
entre parenth\`eses \`a la date d'\'ecriture~--- elles sont toutes massivement utilis\'ees)~:

\begin{tabz}{Xl}
  \headrow \thd{Extension (identifiant \`a coller)} \& \thd{Pour} \\
  \midrule
  \link{ESLint}{https://marketplace.visualstudio.com/items?itemName=dbaeumer.vscode-eslint} \cmd{dbaeumer.vscode-eslint} \& JS/TS : analyse et auto-correction.\\
  \link{Prettier}{https://marketplace.visualstudio.com/items?itemName=esbenp.prettier-vscode} \cmd{esbenp.prettier-vscode} \& Formatage automatique JS/TS/CSS/JSON/MD.\\
  \link{Error Lens}{https://marketplace.visualstudio.com/items?itemName=usernamehw.errorlens} \cmd{usernamehw.errorlens} \& Affiche l'erreur inline, bout de la ligne.\\
  \link{EditorConfig}{https://marketplace.visualstudio.com/items?itemName=EditorConfig.EditorConfig} \cmd{EditorConfig.EditorConfig} \& Respecte indentation/fins de ligne du projet.\\
  \link{GitLens}{https://marketplace.visualstudio.com/items?itemName=eamodio.gitlens} \cmd{eamodio.gitlens} \& Historique ligne par ligne.\\
  \link{PHP Intelephense}{https://marketplace.visualstudio.com/items?itemName=bmewburn.vscode-intelephense-client} \cmd{bmewburn.vscode-intelephense-client} \& PHP : compl\'etion, types, navigation.\\
  \link{PHP Debug}{https://marketplace.visualstudio.com/items?itemName=xdebug.php-debug} \cmd{xdebug.php-debug} \& Xdebug pas \`a pas.\\
  \link{PHP Namespace Resolver}{https://marketplace.visualstudio.com/items?itemName=mehedidrac PHP Namespace Resolver} \& Import automatique des classes.\\
  \link{Extension Pack for Java}{https://marketplace.visualstudio.com/items?itemName=vscjava.vscode-java-pack} \cmd{vscjava.vscode-java-pack} \& Les 6 extensions Java de Red~Hat/Microsoft d'un coup.\\
  \link{Spring Boot Extension Pack}{https://marketplace.visualstudio.com/items?itemName=vmware.vscode-boot-dev-pack} \& Autocompl\'etion \verb|application.properties|, Initializr.\\
  \link{REST Client}{https://marketplace.visualstudio.com/items?itemName=humao.rest-client} \cmd{humao.rest-client} \& Tester une API depuis un fichier \fichier{.http}.\\
  \link{Thunder Client}{https://www.thunderclient.com/} \& Alternative \`a Postman int\'egr\'ee.\\
  \link{Docker}{https://marketplace.visualstudio.com/items?itemName=ms-azuretools.vscode-docker} \& Containers, logs, volumes.\\
  \link{SQLite}{https://marketplace.visualstudio.com/items?itemName=qwtel.sqlite-viewer} \& Explorer une base.\\
  \link{Code Spell Checker}{https://marketplace.visualstudio.com/items?itemName=streetsidesoftware.code-spell-checker} \& \'Evite les fautes dans les noms de variables.\\
  \link{Live Server}{https://marketplace.visualstudio.com/items?itemName=ritwickdey.LiveServer} \& Rechargement automatique d'une page statique.\\
\end{tabz}

R\'eglages que j'active imm\'ediatement (\keys{Ctrl+,}, puis \clic{Open Settings (JSON)})~:

\begin{codejson}
{
  "editor.formatOnSave": true,
  "editor.defaultFormatter": "esbenp.prettier-vscode",
  "editor.tabSize": 2,
  "editor.insertSpaces": true,
  "editor.rulers": [80, 100, 120],
  "editor.bracketPairColorization.enabled": true,
  "editor.guides.bracketPairs": "active",
  "editor.linkedEditing": true,
  "editor.suggestSelection": "first",
  "files.trimTrailingWhitespace": true,
  "files.insertFinalNewline": true,
  "files.eol": "\n",
  "[php]": { "editor.tabSize": 4, "editor.defaultFormatter": "xiaphu.php-build-tools" },
  "[java]": { "editor.tabSize": 4 },
  "javascript.preferences.importModuleSpecifier": "relative",
  "typescript.preferences.importModuleSpecifier": "relative",
  "java.configuration.updateBuildConfiguration": "automatic",
  "java.completion.enabled": true,
  "editor.inlineSuggest.enabled": true,
  "terminal.integrated.defaultProfile.windows": "Ubuntu-24.04 (WSL)"
}
\end{codejson}

\begin{note}
L'identifiant \cmd{xiaphu.php-build-tools} dans le bloc \verb|[php]| est optionnel~:
si tu ne l'as pas install\'e, supprime la ligne \verb|editor.defaultFormatter| du bloc PHP,
sinon VS~Code affichera une erreur de formateur introuvable \`a chaque sauvegarde.
\end{note}

\subsection{Option 2 --- les IDE \og  tout-en-un  \fg{} de r\'ef\'erence}

\begin{tabz}{XlX}
  \headrow \thd{IDE} \& \thd{Prix} \& \thd{Liens} \\
  \midrule
  \link{PhpStorm}{https://www.jetbrains.com/phpstorm/} (JetBrains) \& payant, gratuit pour \'etudiants / open source / enseignant \& le meilleur IDE PHP au monde ; int\`egre Xdebug, Composer, Doctrine, DB~Tools. \\
  \link{IntelliJ~IDEA~Community}{https://www.jetbrains.com/idea/download/} \& gratuit \& l'IDE Java de r\'ef\'erence ; la version Ultimate (payante) ajoute Spring, Jakarta EE, bases de donn\'ees. \\
  \link{VS~Code}{https://code.visualstudio.com/} \& gratuit \& suffisant pour tout ce que contient ce livre. \\
  \link{Eclipse~IDE~for~Java}{https://www.eclipse.org/downloads/} \& gratuit \& encore tr\`es pr\'esent en entreprise et en fac. \\
  \link{Apache NetBeans}{https://netbeans.apache.org/download/index.html} \& gratuit \& excellent avec Java, Maven, GlassFish. \\
  \link{Zed}{https://zed.dev/} \& gratuit \& \'editeur Rust ultra-rapide, multi-langage, en pleine croissance. \\
  \link{Vim/Neovim}{https://neovim.io/} \& gratuit \& investissement long, rendement infini. \\
\end{tabz}

\begin{astuce}
Choisis \textbf{un} \'editeur pour les trois langages et garde-le trois mois. Changer d'outil
 toutes les deux semaines est la forme la plus fr\'equente de procrastination d\'eguis\'ee en
 productivit\'e. Si ton budget le permet et que tu fais beaucoup de PHP, PhpStorm + IntelliJ
 Community est le duo qui fait gagner le plus de temps.
\end{astuce}

\section{Bases de donn\'ees et clients}

\subsection{La m\'ethode Docker (recommand\'ee)}

Tu ne pollues pas ta machine, tu lances et tu jettes une base \`a volont\'e.

\begin{terminal}
# MySQL 8 utilise par le chapitre PHP
docker run -d --name mysql-dev \
  -e MYSQL_ROOT_PASSWORD=secret \
  -e MYSQL_DATABASE=blog \
  -e MYSQL_USER=nour -e MYSQL_PASSWORD=nourpwd \
  -p 3306:3306 -v mysql-data:/var/lib/mysql \
  mysql:8.4 --character-set-server=utf8mb4 --collation-server=utf8mb4_unicode_ci

# PostgreSQL 16 utilise par les chapitres Java
docker run -d --name pg-dev \
  -e POSTGRES_PASSWORD=secret -e POSTGRES_DB=boutique \
  -p 5432:5432 -v pg-data:/var/lib/postgresql/data \
  postgres:16

# Adminer : un client web en un conteneur (alternative legere a phpMyAdmin)
docker run -d --name adminer -p 8080:8080 adminer:5

# Redis (cache/queues, utilise dans la partie Node)
docker run -d --name redis-dev -p 6379:6379 redis:7-alpine

docker ps          # la liste des conteneurs qui tournent
docker logs -f mysql-dev
docker stop mysql-dev && docker start mysql-dev
\end{terminal}

Clients graphiques~: \link{DBeaver~Community}{https://dbeaver.io/} (gratuit, toutes bases, mon pr\'ef\'er\'e),
\link{TablePlus}{https://tableplus.com/} (magnifique, freemium, macOS/Windows/Linux),
\link{HeidiSQL}{https://www.heidisql.com/} (Windows, gratuit, excellent avec MySQL),
\link{phpMyAdmin}{https://www.phpmyadmin.net/} (web), \link{pgAdmin}{https://www.pgadmin.org/}.

\subsection{Sans Docker}
\begin{itemize}
  \item Windows/macOS~: \link{XAMPP}{https://www.apachefriends.org/fr/index.html} ou
        \link{Laragon}{https://laragon.org/} (Windows, le meilleur choix local PHP+MySQL+nginx).
  \item Linux~: \cmd{sudo apt install apache2 mariadb-server php-mysql} puis
        \cmd{sudo a2enmod rewrite && sudo systemctl restart apache2}.
  \item \link{MySQL~Shell}{https://dev.mysql.com/downloads/shell/} et
        \link{Postgres.app}{https://postgresapp.com/} (macOS, z\'ero configuration).
\end{itemize}

\section{Tester une API : le client HTTP}

\begin{itemize}
  \item \link{Postman}{https://www.postman.com/downloads/} --- le standard, gratuit avec compte.
  \item \link{Insomnia}{https://insomnia.rest/download} --- open source, l\'eger, sans compte.
  \item \link{Bruno}{https://www.usebruno.com/} --- collections stock\'ees dans Git, tr\`es appr\'eci\'e des \'equipes.
  \item \link{Hoppscotch}{https://hoppscotch.io/} --- dans le navigateur, rien \`a installer.
  \item \cmd{curl} --- la ligne de commande~: tu l'utiliseras tous les jours dans ce livre.
  \item L'extension \link{REST~Client}{https://marketplace.visualstudio.com/items?itemName=humao.rest-client}
        de VS~Code~: des requ\^etes \'ecrites dans un fichier \fichier{http/api.http} versionn\'e avec le code.
\end{itemize}

\begin{terminal}
# Exemple type d'appel a une API locale
curl -i http://localhost:8000/api/health
curl -s -X POST http://localhost:8000/api/users \
  -H "Content-Type: application/json" \
  -d '{"name":"Nour","email":"nour@example.com"}' | jq .
\end{terminal}

\section{Le minimum syndical : v\'erifier que tout fonctionne}

\begin{terminal}
node -v ; npm -v
php -v ; php -m | grep -E "pdo_mysql|mbstring|curl|openssl" ; composer -V
java -version ; javac -version ; echo $JAVA_HOME
mvn -v ; gradle -v
git --version ; docker -v ; docker compose version
sqlite3 --version ; mysql --version ; psql --version
\end{terminal}

Si \cmd{mvn} ou \cmd{gradle} manquent~: \cmd{sudo apt install maven} et \cmd{brew install maven},
ou utilise les wrappers \cmd{./mvnw} et \cmd{./gradlew} fournis par chaque projet
(voir \cref{ch:java-build}).

\section{R\'eparation express}

\begin{tabz}{Xl}
  \headrow \thd{Sympt\^ome} \& \thd{Cause probable et rem\`ede} \\
  \midrule
  \texttt{command not found: php} \& Le binaire n'est pas dans le \verb|PATH|. \cmd{which php} / \cmd{where.exe php}. Sous Windows, red\'emarre le terminal apr\`es modification du PATH.\\
  \texttt{java: command not found} \& JAVA\_HOME non d\'efini sous Windows ou \verb|PATH| non rafra\^ichi. Relance le terminal.\\
  \texttt{EACCES: permission denied} (npm) \& N''utilise jamais \cmd{sudo npm}. Change le pr\'efixe~: \cmd{npm config set prefix \textasciitilde/.npm-global}.\\
  \texttt{PHP Startup: Unable to load dynamic library} \& Extension comment\'ee ou mauvaise version dans \fichier{php.ini}~; v\'erifie \cmd{php -i | grep "Loaded Configuration"}.\\
  \texttt{docker: permission denied on socket} \& Session non recharg\'ee apr\`es \cmd{usermod -aG docker}. D\'econnexion/reconnexion, ou \cmd{newgrp docker}.\\
  Port 3306/5432/80 d\'ej\`a utilis\'e \& Un serveur local tourne d\'ej\`a~: \cmd{sudo lsof -i :3306} (Linux/macOS) ou \cmd{netstat -ano | findstr :3306} (Windows), puis change le port mapp\'e dans Docker.\\
  WSL lent \& Projet plac\'e dans \fichier{/mnt/c}. D\'eplace-le dans le syst\`eme de fichiers Linux.\\
\end{tabz}

\begin{exercice}
\begin{enumerate}
  \item Installe les trois runtimes et fais afficher leurs versions. Fais une capture d'\'ecran~:
        tu t'en serviras dans les rapports de bug et sur les forums.
  \item Cr\'ee \verb|~/labo/| (ou \fichier{C:\textbackslash labo}), mets-la sous Git avec
        \cmd{git init} et un \fichier{.gitignore} contenant \verb|node\_modules/|,
        \verb|vendor/|, \verb|target/|, \verb|.env|.
  \item Lance les deux conteneurs MySQL et PostgreSQL, v\'erifie avec \cmd{docker ps},
        puis arr\^ete-les proprement~: \cmd{docker stop mysql-dev pg-dev}.
  \item Configure VS~Code pour que \keys{Ctrl+S} formate, enl\`eve les espaces de fin et ajoute
        un saut de ligne final.
  \item Bonus~: installe le th\`eme \link{Catppuccin}{https://github.com/catppuccin/vscode} ou
        \link{One~Dark~Pro}{https://marketplace.visualstudio.com/items?itemName=binaryify.onedarkpro}
        et la police \link{JetBrains~Mono}{https://www.jetbrains.com/lp/mono/} (gratuite,
        avec ligatures~: \verb|=>| et \verb|!=| deviennent lisibles).
\end{enumerate}
\end{exercice}

\begin{ressources}
\begin{itemize}
  \item MDN --- \emph{Installation of Node.js}~: \url{https://developer.mozilla.org/en-US/docs/Learn/Tools\\\\_and\\\\_testing/Understanding\\\\_client-side\\\\_tools/Overview\#installing\\\\_node.js}
  \item PHP --- \emph{Installation}~: \url{https://www.php.net/manual/en/install.php}
  \item OpenJDK --- \url{https://openjdk.org/install/} --- Adoptium Temurin~: \url{https://adoptium.net/fr/}
  \item Composer --- \url{https://getcomposer.org/download/}
  \item Docker Desktop --- \url{https://docs.docker.com/desktop/install/windows-install/}
  \item WSL --- \url{https://learn.microsoft.com/fr-fr/windows/wsl/install}
\end{itemize}
\end{ressources}

\index{installation}\index{outils}\index{VS Code}
